\chapter{The Geometry of the Active Set}
\label{ch:geometry}

Chapter~\ref{ch:certificate} reduced the QP to a search over working sets. This
chapter studies that search space: what a working set is geometrically, how the
optimality conditions read as a \emph{complementarity} problem, why the
bound-constrained case $C = I$ is structurally privileged, and what degeneracy is
and does. It contains no algorithms. It contains the facts that decide which
algorithms are possible.

\section{Faces, working sets, and the lattice}

The feasible set $\Omega$ is a polyhedron. Its faces are obtained by promoting
some subset of the inequalities to equalities: for
$\Active \subseteq \{m_{\mathrm{eq}}+1,\dots,m\}$,
\[
  F_\Active = \{x \in \Omega : c_i\T x = b_i \text{ for } i \in \Active\}
\]
is a face, possibly empty. These faces, ordered by inclusion, form a lattice, and
the solution of~\eqref{eq:qp} lies in the relative interior of exactly one of
them---namely $F_{\Active(x^\star)}$, where $\Active(x^\star)$ is the active set
of Definition~\ref{def:activeset}.
\index{face lattice}

So the QP is a two-stage problem: choose a face, then minimise a strictly convex
quadratic on the affine hull of that face, which is the subproblem~\eqref{eq:sub}.
The second stage is linear algebra. The first stage has
$2^{m - m_{\mathrm{eq}}}$ candidates, and every method in this book is a way of
visiting a vanishingly small number of them.

Two quantities bound the search. The first is elementary and pessimistic.

\begin{proposition}[Working-set size]
\label{prop:worksize}
If $A = C_{\cdot\Active}$ has full column rank then $|\Active| \le n$. In
particular the solution's working set never exceeds $n$ constraints, however
large $m$ is.
\end{proposition}

\begin{proof}
$A \in \R^{n\times k}$ with rank $k$ forces $k \le n$.
\end{proof}

The second is the observation that makes active-set methods practical: the
solution's active set is usually much smaller than $n$, or---in the
bound-constrained case---much \emph{larger}, in a way that is equally exploitable.
For a long-only portfolio the optimum is a vertex at which most non-negativity
bounds bind: on S\&P~500 data with $n = 494$ assets, roughly $450$ of the $495$
constraints are active at the minimum-variance solution and only $46$ assets are
held. Chapter~\ref{ch:krylov} exploits the small \emph{free} set; the synthetic
test families common in the literature, whose active sets are small, do not
exercise the same code paths, and Chapter~\ref{ch:measurement} argues that this
matters for benchmarking.

\section{Complementarity}

Specialise to bound constraints, $C = I$ and $b = 0$, so the problem is
\begin{equation}
  \min_{x \ge 0} \; f(x) = \tfrac12 x\T A x - b\T x, \qquad A \succ 0,
  \label{eq:bqp}
\end{equation}
where we have switched to the symbols $(A,b)$ that the literature on this case
uses; $A$ here plays the role of $G$ above. Introduce the multiplier
$s \ge 0$ for $x \ge 0$. Stationarity gives $s = Ax - b = \nabla f(x)$, the
\emph{reduced gradient}, and the KKT conditions collapse to
\index{linear complementarity problem}
\begin{equation}
  s = A x - b, \qquad x \ge 0, \qquad s \ge 0, \qquad x_i\,s_i = 0 \ \ \forall i.
  \label{eq:lcp}
\end{equation}

\begin{definition}[Linear complementarity problem]
Given $M \in \R^{n\times n}$ and $q \in \R^n$, the problem
$\mathrm{LCP}(M,q)$ is to find $x$ with $x \ge 0$, $Mx + q \ge 0$, and
$x\T(Mx+q) = 0$.
\end{definition}

So~\eqref{eq:lcp} is $\mathrm{LCP}(A,-b)$, and the bound-constrained QP and this
LCP are the same problem. That identification is not a curiosity: it imports a
theory, and Chapter~\ref{ch:guessing} spends the theory.

Complementarity partitions the indices. Write
\[
  \Free = \{i : x_i > 0\}
  \quad\text{(the \emph{free} set)},
  \qquad
  \Bound = \{i : x_i = 0\}
  \quad\text{(the \emph{bound} set)}.
\]
\index{free set}\index{bound set}
On $\Free$ we have $s_i = 0$, hence $(Ax)_i = b_i$; on $\Bound$ we have
$s_i \ge 0$. A bound variable with $s_i < 0$ would lower $f$ if released to a
small positive value; a free variable with $x_i < 0$ is infeasible and must be
pushed to the bound. Those two violations are precisely the primal and dual steps
of the algorithm in Chapter~\ref{ch:guessing}, and they are visible here, before
any algorithm exists, as the only two ways a candidate partition can be wrong.

Once $\Free$ is fixed the sign constraints drop out and what remains is the SPD
linear system
\begin{equation}
  A_{\Free}\,x_{\Free} = b_{\Free},
  \qquad A_\Free := A_{\Free\Free},
  \label{eq:innerspd}
\end{equation}
which is~\eqref{eq:sub} written for this case. Chapters~\ref{ch:krylov}
and~\ref{ch:operator} are about solving~\eqref{eq:innerspd} well.

\section{$P$-matrices: why bounds are privileged}

Here is the structural fact that separates the bound-constrained case from the
general one, and it is worth stating carefully because a great deal follows from
it.

\begin{definition}[$P$-matrix]
\index{P-matrix@$P$-matrix}
A matrix $M \in \R^{n\times n}$ is a \emph{$P$-matrix} if all its principal minors
are positive.
\end{definition}

\begin{proposition}[SPD implies $P$]
\label{prop:spdp}
If $A \succ 0$ then every principal submatrix $A_{\Free}$ is symmetric positive
definite, hence $\det A_\Free > 0$; so $A$ is a $P$-matrix.
\end{proposition}

\begin{proof}
For nonzero $x_\Free$, extend by zeros to $x \in \R^n$; then
$x_\Free\T A_\Free x_\Free = x\T A x > 0$. Symmetry is inherited. Positive
definiteness gives positive determinant.
\end{proof}

Two consequences, of quite different depth. The shallow one: \emph{every}
candidate free set gives a well-posed inner system~\eqref{eq:innerspd}, so an
algorithm may guess wildly and still have something to solve. The deep one:
$\mathrm{LCP}(M,q)$ with $M$ a $P$-matrix has a \emph{unique} solution for every
$q$, and that solution is reachable by principal pivoting in finitely many steps
\cite{murty1988,judice1994}. This is the hypothesis behind the finite-termination
theorem of Chapter~\ref{ch:guessing}.

Now the general case. With constraints $C\T x \ge b$ and a guessed working set
$\Active$, the system to be solved is
\begin{equation}
  \bigl(C_\Active\T G^{-1} C_\Active\bigr)\lambda_\Active
   = b_\Active - C_\Active\T x_u ,
  \label{eq:generalsystem}
\end{equation}
whose coefficient matrix is the dual Hessian of Definition~\ref{def:dualhessian}.
Compare the two cases:

\begin{center}
\begin{tabular}{lll}
\toprule
 & Bound constraints ($C = I$) & General constraints \\
\midrule
Working-set system & $(G^{-1})_{\Active\Active}$ & $C_\Active\T G^{-1} C_\Active$ \\
Structure & principal submatrix of $G^{-1}$ & Gram matrix of $C_\Active$ in $G^{-1}$ \\
Positive definite when & \emph{always} & $C_\Active$ has full column rank \\
That is a property of & the data & the \emph{guess} \\
\bottomrule
\end{tabular}
\end{center}

\index{dual Hessian!rank condition}
When $C = I$, equation~\eqref{eq:generalsystem} reads
$(G^{-1})_{\Active\Active}\lambda_\Active = \cdots$, a principal submatrix of
$G^{-1} \succ 0$, hence positive definite whatever $\Active$ is. When $C$ is
general, $C_\Active\T G^{-1} C_\Active$ is positive definite exactly when
$C_\Active$ has full column rank---and a guessed set may name $k > n$ constraints,
or $k \le n$ linearly dependent ones. There is no $P$-matrix to appeal to, and
correspondingly no theorem to inherit.

This is the single most important structural difference in the book. It is not a
gap in an argument or a missing anti-cycling rule; it is a property of the two
constraint classes. Chapter~\ref{ch:guessing} measures how often it bites (rarely
on well-posed random data, and in $93\%$ of instances once columns of $C$ are
duplicated), and Chapter~\ref{ch:walking} shows the elegant compensation: the dual
method's \emph{step rules} make linear dependence unreachable, so it needs no rank
test at all.

\section{Degeneracy}

\begin{definition}[Primal and dual degeneracy]
\index{degeneracy}
A feasible point $x$ is \emph{primal degenerate} if more constraints are active at
$x$ than the working set holds---equivalently, if a constraint outside $\Active$
has zero slack. A KKT point is \emph{dual degenerate} (or fails \emph{strict
complementarity}) if some $i \in \Active(x^\star)$ has $\lambda_i = 0$.
\end{definition}

Both are non-generic: each is a codimension-one condition on the data, so a
random instance misses them with probability one. Both nevertheless occur
constantly in practice, because real data is not random. Duplicated assets,
tied expected returns, a constraint entered twice by a model-assembly script, a
group budget that happens to coincide with a sum of individual bounds---each
produces exact ties.

The consequences are specific and worth separating.

\paragraph{Primal degeneracy stalls ascent arguments.} A method that proves
termination by showing the objective strictly increases at each step is defeated
by a step of length zero, and a zero step is exactly what primal degeneracy
permits: a constraint with zero slack can be added to the working set without
moving $x$ at all. The objective does not change, the ascent argument goes silent,
and a cycle is not excluded. This is the hypothesis in
Proposition~\ref{prop:ascent} of Chapter~\ref{ch:walking}.

\paragraph{Dual degeneracy blurs the identification of the active set.} If
$\lambda_i = 0$ for an active $i$, then $i$ may be in or out of the working set at
the solution without violating anything, so the working set at the optimum is not
unique. Methods that promise to ``identify the optimal face finitely'' generally
assume strict complementarity, and Chapter~\ref{ch:krylov} contrasts a family that
does with the method it develops, which does not.

\paragraph{Ties need a rule.} When several events would occur at the same
parameter value, choosing among them arbitrarily can cycle. The remedy is
inherited from linear programming: \emph{Bland's rule}, which breaks ties by
lowest index. It makes the walk deterministic and finite. Murty's least-index
rule for LCPs \cite{murty1988} is the same device, and it is what
Chapter~\ref{ch:guessing}'s termination theorem rests on. Lexicographic
perturbation is the other classical remedy.
\index{Bland's rule}

\begin{remark}[Degeneracy in floating point]
\label{rem:degen-float}
In finite precision the question changes character, because a slack of exactly
zero is not what one observes. What one observes is a slack of $10^{-16}$ or
$-10^{-16}$, and whether a constraint is ``active'' becomes a decision made by a
threshold. This is why Chapter~\ref{ch:numerics} treats degeneracy as a numerical
subject as much as a combinatorial one, and why the two tolerances---the one
deciding what to try next and the one deciding what to return---have very
different jobs and very different delicacy.
\end{remark}

\section{How many steps? The combinatorial question}

An active-set method changes the working set by one index per step, so its step
count is a walk length in the face lattice. Two facts frame what is known.

Empirically, the count is near-linear. A dual method on the QP typically takes
$O(n)$ outer steps; a parametric trace of a frontier typically has $O(n)$
breakpoints; a block-pivoting method converges in a handful of repairs almost
independently of $n$. These are the observations that make the subject practical.

In the worst case, it is not. Mairal and Yu \cite{mairal2012} construct LASSO
paths whose length is \emph{exponential} in the dimension, and the corresponding
bound for multi-parametric quadratic programs---exponentially many critical
regions---has been known in control since \cite{tondel2003}. The simplex method's
history is the obvious analogue: near-linear in practice, exponential on
Klee--Minty examples, and the gap between the two is still not fully explained.

\emph{What separates the benign average from the adversarial worst case is open},
in the shared form that covers all the instances of Part~III. It is the most
interesting unsolved problem this book touches, and it is stated here rather than
in a footnote because a graduate reader should know which parts of a subject are
settled.

\begin{notes}
Cottle, Pang and Stone \cite{cottle1992} is the standard reference for linear
complementarity, Murty \cite{murty1988} for principal pivoting and the least-index
rule, and Júdice and Pires \cite{judice1994} for the block-pivoting construction
that Chapter~\ref{ch:guessing} uses. The $P$-matrix characterisation of LCP
uniqueness is in both.

The $93\%$ figure, and the claim that the rank obstruction never occurs on
$600$ well-posed random instances across four constraint families, are measured in
the Goldfarb--Idnani note \cite{schmelzer2026quadprog}. The observation that a long-only optimum is a vertex holding
few names is from \cite{schmelzer2026matrixfree}.
\end{notes}

\begin{exercises}

\begin{exercise}
Show that the face $F_\Active$ is non-empty for at most $\binom{m}{n}$ of the
$2^m$ index sets when $\Omega$ is bounded and in general position. Why does this
not make enumeration practical?
\end{exercise}

\begin{exercise}
Prove that if $M$ is a $P$-matrix then $\mathrm{LCP}(M,q)$ has at most one
solution. (Hint: if $x \ne y$ both solve it, consider the index $i$ maximising
$|x_i - y_i|$ and the sign of $(M(x-y))_i(x-y)_i$.)
\end{exercise}

\begin{exercise}
Give a $2 \times 2$ example of a $P$-matrix that is not symmetric and not positive
definite. Deduce that Proposition~\ref{prop:spdp} is a strict implication.
\end{exercise}

\begin{exercise}
Let $C$ have two identical columns, $c_1 = c_2$, with $b_1 = b_2$. Show the
feasible set is unchanged by deleting one of them, but that a guessed working set
containing both makes~\eqref{eq:generalsystem} singular. This is the mechanism
behind the $93\%$ figure of Chapter~\ref{ch:guessing}.
\end{exercise}

\begin{exercise}
Construct a two-variable QP whose solution is primal degenerate: three constraints
passing through the optimum in $\R^2$. Verify that two different working sets both
satisfy the KKT conditions there, with different multipliers.
\end{exercise}

\begin{exercise}
For the bound-constrained problem~\eqref{eq:bqp}, show that $x$ solves it if and
only if $x = \max(x - \tau(Ax - b),\,0)$ for any $\tau > 0$, where the maximum is
componentwise. (This fixed-point characterisation is the basis of the projected
gradient method of Chapter~\ref{ch:krylov}.)
\end{exercise}

\begin{exercise}[Implementation]
Generate random long-only minimum-variance instances with $n = 20, 50, 100, 200$
from a factor model, solve each with any QP solver, and plot the number of assets
held (the free set size) against $n$. Confirm that it grows far more slowly than
$n$. Now repeat with $\Sig = I$ and explain the difference.
\end{exercise}

\end{exercises}
